\section{Summary}
\begin{frame}{}
    \LARGE Policy Optimization: \textbf{Summary}
\end{frame}

\begin{frame}{The Algorithm Family Tree}
\begin{center}
REINFORCE $\;\rightarrow\;$ +baseline $\;\rightarrow\;$ +critic (A2C) $\;\rightarrow\;$ +GAE $\;\rightarrow\;$ +trust region (TRPO) $\;\rightarrow\;$ +clipping (PPO) $\;\rightarrow\;$ $-$critic (GRPO)
\end{center}

\vspace{0.5em}

\begin{table}[]
    \centering
    \renewcommand{\arraystretch}{1.8}
    \small
    \begin{tabular}{lcccc}
        \hline
        \textbf{Algorithm} & \textbf{On-policy?} & \textbf{Variance} & \textbf{Sample eff.} & \textbf{Complexity} \\
        \hline
        REINFORCE        & Yes & High   & Low    & Simple \\
        A2C              & Yes & Medium & Medium & Moderate \\
        TRPO             & Yes & Medium & Medium & High \\
        PPO              & Yes & Medium & Medium & Low \\
        GRPO             & Yes & Medium & Medium & Low (no critic) \\
        \hline
    \end{tabular}
\end{table}
\end{frame}

\begin{frame}{Summary (1/2): Getting a Good Advantage Estimate}
\begin{itemize}
    \item \textbf{TD learning} --- update a prediction using a later, better-informed one. The TD error $\delta_t = r_t + \gamma V_\phi(s_{t+1}) - V_\phi(s_t)$ is the same bracket you saw in Q-learning and SARSA; the critic is trained by minimising its \emph{square}
    \item A \textbf{critic} $V_\phi$ turns Day 3's baseline into something learnable, and $\delta_t$ doubles as an online, low-variance estimate of the advantage --- one number, two jobs
    \item \textbf{GAE} replaces the crude 1-step $\delta_t$ with a $\lambda$-weighted average over all lookahead lengths, computed backwards in one pass: $\hat{A}_t = \delta_t + \gamma\lambda \hat{A}_{t+1}$. Read $\lambda$ as ``how much do I trust my critic''
    \item $\hat{A}_t$ is what \emph{every} algorithm from here on multiplies the log-probability gradient by
    \item \textbf{A2C} parallelises data collection with $N$ synchronous actors; A3C uses asynchronous workers (empirically A2C $\geq$ A3C)
\end{itemize}
\end{frame}

\begin{frame}{Summary (2/2): Taking a Safe Step}
\begin{itemize}
    \item A large step in \emph{weights} can be a catastrophic step in \emph{behaviour} --- and a collapsed policy collects bad data, so it rarely recovers
    \item To take a step safely we must score a candidate policy \textbf{without running it}. The \textbf{surrogate objective} $\mathcal{L}(\theta) = \mathbb{E}_t[r_t(\theta)\,\hat{A}_t]$ does this by re-counting already-collected samples by the ratio $r_t(\theta) = \pi_\theta(a_t|s_t)/\pi_{\theta_\text{old}}(a_t|s_t)$
    \item That score is only trustworthy while the two policies stay close --- hence every method below is ``maximize $\mathcal{L}$, but don't move far''
    \item \textbf{TRPO} enforces this with a hard KL constraint. Correct, but it needs the Fisher matrix: second-order, expensive, painful to implement
    \item \textbf{PPO-Clip} gets the same effect first-order: clip $r_t(\theta)$ to $[1{-}\varepsilon, 1{+}\varepsilon]$ and take the $\min$ with the unclipped term. One combined loss = clipped objective $-\,c_1\times$ critic error $+\,c_2\times$ entropy
    \item \textbf{GRPO} drops the critic entirely, using a group-mean baseline --- now the default for LLM RL (DeepSeek-R1)
\end{itemize}
\end{frame}

\begin{frame}{Limitations of Actor-Critic \& PPO}
\begin{itemize}
    \item \textbf{Function approximation bias:} the TD residual $\delta_t$ is only an unbiased advantage estimate if $V_\phi = V^{\pi}$. An inaccurate critic introduces bias into the policy gradient
    \item \textbf{Two-network instability:} if the critic lags behind the actor (or vice versa), the feedback loop can diverge --- both networks must be trained carefully together (GRPO sidesteps this by dropping the critic entirely)
    \item \textbf{Hyperparameter sensitivity:} PPO performance is sensitive to $\varepsilon$ (clip range), $\lambda$ (GAE), number of epochs per rollout, and entropy coefficient $c_2$
    \item \textbf{On-policy sample inefficiency:} all methods in this lecture discard experience after each update. In complex environments this can require billions of environment steps
    \item \textbf{Dense-reward assumption:} variance reduction via advantage estimation works well when rewards are frequent. Sparse-reward tasks remain challenging for all on-policy methods
\end{itemize}
\end{frame}
